%% ============================================================
%%  chapters/chapter1.tex - Introduction
%% ============================================================
\chapter{Introduction}
\label{chap:introduction}

\section{Background}
\label{sec:background}

Programming assessments are increasingly delivered through online learning management systems, browser-based coding environments, and automated grading platforms. These systems are convenient and scalable, but they create a persistent academic integrity problem: the final submitted code alone does not prove that the enrolled student personally produced the work during the assessment session. Password login verifies access at the beginning of a session, but it does not provide continuing evidence of authorship after login. Static plagiarism detectors compare final code similarity, but they cannot reliably explain whether a solution was typed gradually, pasted from an external source, written by another person, or produced through genuine struggle and correction.

Keystroke dynamics provides a promising answer because it treats typing rhythm as a behavioural biometric. Features such as dwell time, flight time, inter-key latency, and rhythm consistency can be captured without extra hardware while the student naturally writes code. Early work by Monrose and Rubin showed that typing patterns can be used for biometric authentication \cite{monrose2000keystroke}. Large-scale studies later confirmed that human typing contains stable personal structure while also varying by task and context \cite{dhakal2018observations}. More recent deep learning work, especially TypeNet, demonstrated that sequence models can transform keystroke timings into compact biometric embeddings suitable for verification and identification \cite{acien2022typenet}.

At the same time, programming education research has shown that interaction traces are not only security evidence. Keystrokes, deletions, pauses, and edit bursts can reveal how students approach a problem, where they struggle, and whether they recover through debugging or abandon the task. Studies on self-regulated learning and programming trace logs show that student process data can support formative intervention, not only summative judgement \cite{arakawa2021situ,dong2021trace}.

This dissertation focuses on this intersection. It presents BLAIM, a Behaviour Learning Analytics and Identity Monitoring module for continuous keystroke authentication and behaviour analysis using keystroke data in programming assessments. The module is implemented inside the GradeLoop Core project and evaluated using recorded enrollment and testing data collected through the BLAIM evaluator environment with real users.

\section{Research Gap}
\label{sec:research-gap}

The literature and current educational technology practice reveal five important gaps.

\textbf{First, authentication and learning analytics are usually separated.} Keystroke biometric systems often focus on accept/reject identity decisions, while learning analytics systems focus on student progress, struggle, and intervention. Continuous authentication research has advanced substantially for free-text typing and online examination scenarios \cite{alshehri2018iterative,lu2020continuous,yang2021tkca,senarath2023reevaluating}, and programming education research has separately shown that coding keystrokes and IDE traces can reveal learning process information \cite{edwards2020keystroke,edwards2023csedm,leinonen2022ide}. However, the literature provides much less work on continuous identity monitoring specifically while students are writing source code. This separation is central to the novelty of BLAIM: in online programming assessments, the same keystroke stream can support both authorship assurance and process-aware learning analytics, but existing systems usually duplicate collection pipelines or ignore one side of the evidence.

\textbf{Second, many authentication outputs are difficult for instructors to interpret.} A similarity score or risk score is useful for a system, but instructors need context: Was the session rejected because the rhythm changed, because too little data was available, because another user matched more closely, or because the process looked like bulk paste behaviour?

\textbf{Third, continuous assessment requires more than login-time verification.} A student may begin a session honestly and later receive help, switch typists, or paste generated code. Login authentication cannot detect these mid-session changes. Continuous monitoring is needed, but it must be designed carefully to avoid intrusive surveillance.

\textbf{Fourth, educational deployment requires clear service boundaries and operational evidence.} Many research prototypes report algorithms but do not provide deployable APIs, data models, archive strategies, dashboards, and integration points for real LMS ecosystems.

\textbf{Fifth, ethical and fairness concerns are not optional.} Keystroke data is biometric and behavioural data. It can be affected by keyboard device, motor ability, stress, injury, language familiarity, and accessibility tools. A responsible system must be positioned as decision support with consent, retention limits, human review, and calibration.

The gap addressed in this work is therefore not simply ``build another keystroke classifier''. The research gap is the lack of an integrated, interpretable, deployable module that performs continuous keystroke identity monitoring during live coding and connects that identity evidence with behaviour analysis for programming assessment workflows.

\section{Research Problem}
\label{sec:research-problem}

The research problem is stated as follows:

\begin{quote}
\textit{How can keystroke data captured during programming assessments be used to continuously authenticate student identity while also producing interpretable behaviour analysis that supports academic integrity review and formative learning support?}
\end{quote}

This problem contains three sub-problems:

\begin{enumerate}
  \item How can raw browser keystroke events be transformed into biometric features suitable for continuous authentication?
  \item How can the same keystroke stream be analysed for meaningful behaviour indicators such as pauses, deletions, paste behaviour, friction points, and engagement?
  \item How can these signals be operationalized in a microservice architecture that integrates with GradeLoop and provides instructor-facing evidence rather than opaque model output?
\end{enumerate}

\section{Aim and Objectives}
\label{sec:aim-objectives}

The aim of this dissertation is to design, implement, and evaluate a continuous keystroke authentication and behaviour analysis module for programming assessments.

The research objectives are:

\begin{enumerate}
  \item To analyse existing research on keystroke dynamics, continuous authentication, programming trace analytics, and online assessment integrity.
  \item To design a system architecture that captures keystroke events once and reuses them for biometric verification, identification, monitoring, archiving, and behaviour analysis.
  \item To implement the module as a GradeLoop-compatible microservice using FastAPI, PostgreSQL, Redis, TypeNet-style embedding inference, and evaluator dashboards.
  \item To evaluate the implementation using recorded enrollment and test sessions, including genuine and impostor trials, and to calculate practical metrics such as genuine accept rate, false rejection, false acceptance, similarity, and risk.
  \item To identify ethical, security, social, operational, and commercialization considerations for deploying keystroke-based assessment integrity tools in educational institutions.
\end{enumerate}

\section{Novelty and Contribution}
\label{sec:novelty}

The novelty of this research is the integration of two normally separate concerns: continuous identity assurance and process-aware behaviour analysis during source-code production. Existing keystroke-authentication studies show that users can be authenticated continuously from typing behaviour, but they are generally not designed around the distinctive structure of programming sessions. Existing programming-trace studies show that coding events can explain learning behaviour, but they usually do not provide biometric identity continuity. BLAIM occupies this underexplored intersection. The contribution is practical and applied: BLAIM does not only compute a biometric match. It also records how the coding session unfolded and presents evidence that can help instructors decide whether a case requires academic integrity review, formative support, or normal acceptance.

The main contributions are:

\begin{enumerate}
  \item A unified event pipeline that captures keystroke data once and supports enrollment, verification, 1:N identification, continuous monitoring, archiving, playback, and analytics.
  \item A TypeNet-inspired biometric pipeline using 70-event sequences, five temporal/key features, and 128-dimensional embeddings for student templates.
  \item A behaviour analysis model that converts keystroke events into interpretable indicators: deletion rate, pause frequency, long pauses, burst typing, rhythm consistency, friction points, authenticity, engagement, active problem solving, and learning depth.
  \item A deployable microservice implementation with REST endpoints, WebSocket monitoring, PostgreSQL persistence, Redis buffering, optional RabbitMQ events, and optional Gemini-assisted analysis.
  \item An evaluation evidence pack based on real user testing data: 32 enrollment records, 31 test records, 34 archives, 12 timeline events, and labelled genuine/impostor outcomes.
\end{enumerate}

Table~\ref{tab:novelty-gap} summarizes how the identified gap is converted into a creative solution.

\begin{table}[H]
  \centering
  \caption{Gap-to-Contribution Summary}
  \label{tab:novelty-gap}
  \begin{singlespace}
  \begin{tabularx}{\textwidth}{@{}p{3.6cm}p{4.6cm}p{5.0cm}@{}}
    \toprule
    \textbf{Existing approach} & \textbf{Observed gap} & \textbf{BLAIM contribution} \\
    \midrule
    Webcam proctoring & Observes the student and environment but does not explain how code was produced & Uses keystroke rhythm and editing behaviour as continuous authorship and process evidence \\
    Static plagiarism checking & Compares final code but misses the construction process & Adds session-level evidence: typing, pauses, deletions, paste behaviour, and archives \\
    Login-time authentication & Verifies access only at the beginning of a session & Performs repeated verification and monitoring during the assessment session \\
    Isolated biometric tools & Provide identity scores without pedagogical interpretation & Combines identity confidence with behaviour analysis for instructor review \\
    Research prototypes & Often lack LMS-ready deployment paths & Implements a GradeLoop-compatible microservice with APIs, database persistence, and evaluator dashboards \\
    \bottomrule
  \end{tabularx}
  \end{singlespace}
\end{table}

\section{Scope}
\label{sec:scope}

The scope of this work is continuous keystroke authentication and behaviour analysis for programming assessments in the GradeLoop Core ecosystem. The system captures keyboard timing and key metadata from coding sessions, creates enrollment templates, verifies new sessions, identifies likely users, monitors active sessions, and archives completed sessions for review.

The following are outside the primary scope:

\begin{itemize}
  \item Webcam proctoring, audio proctoring, and face recognition.
  \item Legal disciplinary automation without instructor review.
  \item Large-scale demographic fairness validation across many institutions.
  \item Semantic code plagiarism detection, although BLAIM can complement such systems.
  \item Production-grade encryption of stored biometric templates, which is recommended for deployment but not fully implemented in the prototype database.
\end{itemize}

\section{Methodological Overview}
\label{sec:method-overview}

The study follows a design science and applied software engineering methodology. First, literature was reviewed to identify the academic and practical gap. Second, functional and non-functional requirements were derived for an LMS-integrated keystroke service. Third, the BLAIM module was designed and implemented within the GradeLoop architecture. Fourth, the module was evaluated using recorded real-user testing and evaluator database records. Finally, the results were discussed with attention to accuracy, usability, limitations, ethics, and commercialization.

\section{Dissertation Structure}
\label{sec:structure}

Chapter 2 reviews related work and compares BLAIM with existing approaches. Chapter 3 explains the methodology, architecture, implementation, technology choices, requirements, testing strategy, and commercialization plan. Chapter 4 presents evaluation results, metrics, findings, discussion, limitations, ethics, and student contribution. Chapter 5 concludes the dissertation and recommends future work.

\clearpage
